\chapter{BPE Tokenization}
\label{ch:bpe}

In the BabyGPT tutorial we used a character-level tokenizer: each of the 65 unique characters in TinyShakespeare mapped to one integer.  That approach is beautifully simple, but it does not scale.  A single English word like ``unsupervised'' becomes 12 tokens, inflating every sequence and forcing the model to spend capacity learning spelling.  Modern language models use \textbf{Byte Pair Encoding} (BPE), which compresses text into a vocabulary of subword tokens --- typically 32\,000 to 100\,000 entries.

In this chapter we build a BPE tokenizer from scratch, train it on real data, and wrap it with \code{tiktoken} for efficient encoding at inference time.  Everything lands in \code{microgpt/tokenizer.py}.


\section{Why Character-Level Doesn't Scale}

Consider the sentence ``The cat sat on the mat.''  A character-level tokenizer produces 27 tokens (one per character including spaces and the period).  A BPE tokenizer with a 32K vocabulary compresses this to about 7 tokens.

That ${\sim}4\times$ reduction matters:
\begin{itemize}
  \item \textbf{Context window.}  A 1\,024-token context holds $\sim$4\,000 characters under BPE versus 1\,024 characters at the character level.  The model sees more text at once.
  \item \textbf{Compute.}  Self-attention is $O(T^2)$ in sequence length $T$.  Reducing $T$ by $4\times$ cuts attention cost by $16\times$.
  \item \textbf{Semantics.}  BPE tokens roughly correspond to words or word pieces, giving each token more meaning than a single character.
\end{itemize}

The tradeoff is tokenizer complexity.  BPE requires a trained merge table and a nontrivial encoding algorithm --- but the payoff in model quality and efficiency is enormous.


\section{Byte Pair Encoding: The Algorithm}

BPE was originally a data compression algorithm (Gage, 1994).  Sennrich et al.\ (2016) adapted it for NLP.  GPT-2 (Radford et al., 2019) introduced the \textbf{byte-level} variant we implement here.

The basic tutorial already walked through the BPE algorithm on words.  Here we focus on the byte-level version that GPT-2 uses.

\subsection{Byte-level base vocabulary}

Instead of starting with Unicode characters, the base vocabulary is all 256 raw byte values (\code{0x00} through \code{0xFF}).  Every possible input --- English, Chinese, emoji, even binary data --- can be represented as a sequence of bytes, so there is \textbf{no unknown-token problem}.

The training corpus is first encoded to UTF-8 bytes, then BPE merges operate on those byte sequences.  For example, the character \texttt{\"u} (U+00FC) becomes the two-byte sequence \texttt{[0xC3, 0xBC]} in UTF-8.  BPE might later merge those two bytes into a single token if the pair is frequent enough.

\subsection{Pre-tokenization}

Before BPE sees the bytes, we split the text into coarse chunks using a regex pattern.  This prevents merges from crossing word boundaries --- we do not want BPE to merge the space at the end of ``the'' with the ``c'' at the start of ``cat'' into a single token.

GPT-2 and GPT-4 use carefully designed split patterns.  Here is the one we adopt (matching GPT-4's pattern with a small tweak for number grouping):

\concepttag
\begin{lstlisting}[language={},basicstyle=\ttfamily\footnotesize]
'(?i:[sdmt]|ll|ve|re)       # contractions: 's 't 'd 'm 'll 've 're
|[^\r\n\p{L}\p{N}]?+\p{L}+ # optional punct/symbol then letters
|\p{N}{1,2}                 # 1-2 digit number groups
| ?[^\s\p{L}\p{N}]++[\r\n]* # punctuation runs + trailing newlines
|\s*[\r\n]                  # standalone newlines
|\s+(?!\S)                  # trailing whitespace
|\s+                        # other whitespace
\end{lstlisting}

Each alternation captures a specific kind of chunk.  The \code{\textbackslash p\{L\}} and \code{\textbackslash p\{N\}} classes match Unicode letters and numbers respectively --- the stdlib \code{re} module does not support these, so we use the \code{regex} library.

\tip{The pre-tokenization regex is not part of the BPE algorithm itself.  It is a \emph{preprocessing step} that constrains which byte sequences BPE is allowed to merge.  Different models use different patterns --- GPT-2, GPT-3, GPT-4, and LLaMA all have slightly different regexes.  The BPE training and encoding algorithms are the same regardless.}

\subsection{The merge loop}

After pre-tokenization, each chunk is a list of byte values (integers 0--255).  The merge loop is identical to what we saw in the basic tutorial, just operating on bytes instead of characters:

\begin{enumerate}
  \item \textbf{Count pairs.}  For every adjacent pair of tokens across all chunks, count the total frequency.
  \item \textbf{Merge the top pair.}  Find the most frequent pair.  Assign it a new token ID (starting from 256).  Replace every occurrence of that pair across all chunks with the new ID.
  \item \textbf{Repeat} until the desired vocabulary size is reached.
\end{enumerate}

Each merge reduces the total number of tokens (better compression) and adds one entry to the vocabulary.  If we want a vocabulary of 32\,768 tokens, we perform $32\text{,}768 - 256 - n_{\text{special}} = 32\text{,}511$ merges (subtracting the 256 byte tokens and any special tokens).


\section{Implementation}

Here is the complete tokenizer.  We will walk through each piece.

\subsection{Constants}

\filetag{microgpt/tokenizer.py}
\begin{lstlisting}
"""BPE tokenizer: train from scratch, encode/decode via tiktoken.

This module implements byte-level Byte Pair Encoding (BPE) as used by GPT-2
and later models.  The training algorithm is pure Python so every step is
visible.  Once trained, the merge table is handed to tiktoken for fast
encoding at inference time.

Typical usage:
    # Train on a corpus
    tokenizer = BPETokenizer.train(corpus_text, vocab_size=32768)
    tokenizer.save("tokenizer_dir")

    # Later: load and use
    tokenizer = BPETokenizer.load("tokenizer_dir")
    ids = tokenizer.encode("Hello, world!")
    text = tokenizer.decode(ids)
"""

import os
import pickle
import regex  # 're' doesn't support \p{L} Unicode categories; use 'regex'

import tiktoken


# ---------------------------------------------------------------------------
# Pre-tokenization split pattern (GPT-4 style)
# ---------------------------------------------------------------------------
# Before BPE sees the text, we split it into coarse chunks so that merges
# never cross word or number boundaries.  Each group in the alternation
# captures a specific kind of token:
#
#   'contractions    --- 's, 't, 'd, 'm, 'll, 've, 're (case-insensitive)
#   letter words     --- optional leading non-letter/non-digit, then letters
#   numbers          --- 1-2 digit groups (keeps numbers short)
#   punctuation runs --- non-letter/non-digit/non-space followed by newlines
#   lone newlines    --- whitespace-only lines
#   trailing space   --- spaces not followed by non-space (eaten separately)
#   other spaces     --- remaining whitespace
#
# The \p{L} and \p{N} classes are Unicode-aware (handled by the 'regex'
# library, not stdlib 're').

SPLIT_PATTERN = (
    r"'(?i:[sdmt]|ll|ve|re)"
    r"|[^\r\n\p{L}\p{N}]?+\p{L}+"
    r"|\p{N}{1,2}"
    r"| ?[^\s\p{L}\p{N}]++[\r\n]*"
    r"|\s*[\r\n]"
    r"|\s+(?!\S)"
    r"|\s+"
)

# Special tokens prepended/appended to documents during training.
SPECIAL_TOKENS = ["<|bos|>"]
\end{lstlisting}

The \code{SPLIT\_PATTERN} is a single regex with seven alternation branches, each matching a different type of text chunk.  The \code{SPECIAL\_TOKENS} list defines tokens that are never produced by the merge algorithm --- they are reserved for structural markers.  For pretraining we only need \code{<|bos|>} (beginning of sequence), which we prepend to every document.

\subsection{Pair counting and merging}

The two core helper functions are \code{\_get\_pair\_counts} (count every adjacent pair) and \code{\_merge\_pair} (replace all occurrences of a pair with a new token):

\filetag{microgpt/tokenizer.py}
\begin{lstlisting}
def _get_pair_counts(
    chunk_ids_list: list[list[int]],
) -> dict[tuple[int, int], int]:
    """Count every adjacent pair across all chunks."""
    counts: dict[tuple[int, int], int] = {}
    for ids in chunk_ids_list:
        for i in range(len(ids) - 1):
            pair = (ids[i], ids[i + 1])
            counts[pair] = counts.get(pair, 0) + 1
    return counts


def _merge_pair(
    ids: list[int],
    pair: tuple[int, int],
    new_id: int,
) -> list[int]:
    """Replace every occurrence of *pair* in *ids* with *new_id*."""
    out: list[int] = []
    i = 0
    while i < len(ids):
        if i < len(ids) - 1 and ids[i] == pair[0] and ids[i + 1] == pair[1]:
            out.append(new_id)
            i += 2
        else:
            out.append(ids[i])
            i += 1
    return out
\end{lstlisting}

These are the same functions from the basic tutorial's BPE pseudocode, adapted to work on integer IDs (bytes and merge tokens) instead of character strings.

\subsection{The training function}

\filetag{microgpt/tokenizer.py}
\begin{lstlisting}
def train_bpe(
    text: str,
    vocab_size: int,
    *,
    verbose: bool = True,
) -> list[tuple[tuple[int, int], int]]:
    """Train a byte-level BPE tokenizer on *text*.

    Returns a list of merges: ``[(pair, new_id), ...]`` in priority order.
    The base vocabulary is the 256 raw byte values (0-255), so the first
    merge produces token 256, the second 257, and so on.

    Args:
        text: Training corpus (plain string).
        vocab_size: Target vocabulary size including the 256 byte tokens.
        verbose: Print progress every 1000 merges.
    """
    assert vocab_size >= 256, "vocab_size must be >= 256 (the byte alphabet)"
    num_merges = vocab_size - 256 - len(SPECIAL_TOKENS)
    assert num_merges >= 0, "vocab_size too small to fit byte alphabet + special tokens"

    # Step 1: Pre-tokenize the text into chunks using the split pattern.
    # This prevents merges from crossing word/number boundaries.
    pat = regex.compile(SPLIT_PATTERN)
    chunks = pat.findall(text)

    # Step 2: Convert each chunk to a list of raw byte values.
    chunk_ids_list: list[list[int]] = [
        list(chunk.encode("utf-8")) for chunk in chunks
    ]

    # Step 3: Iteratively merge the most frequent adjacent pair.
    merges: list[tuple[tuple[int, int], int]] = []
    for i in range(num_merges):
        pair_counts = _get_pair_counts(chunk_ids_list)
        if not pair_counts:
            break  # nothing left to merge
        best_pair = max(pair_counts, key=pair_counts.get)
        new_id = 256 + i
        # Apply the merge to every chunk
        chunk_ids_list = [
            _merge_pair(ids, best_pair, new_id) for ids in chunk_ids_list
        ]
        merges.append((best_pair, new_id))
        if verbose and (i + 1) % 1000 == 0:
            print(
                f"  merge {i + 1}/{num_merges}: "
                f"{best_pair} -> {new_id} "
                f"(freq {pair_counts[best_pair]})"
            )

    if verbose:
        print(f"Trained {len(merges)} merges (vocab_size={256 + len(merges) + len(SPECIAL_TOKENS)})")

    return merges
\end{lstlisting}

Walk through the logic:

\begin{enumerate}
  \item \textbf{Pre-tokenize}: split the corpus into chunks with the regex pattern.  Each chunk is a word, number, or punctuation run.
  \item \textbf{Byte-encode}: convert each chunk from a Python string to a list of UTF-8 byte values (integers 0--255).
  \item \textbf{Merge loop}: at each iteration, count all adjacent pairs across all chunks, find the most frequent one, assign it the next available ID ($256 + i$), and replace every occurrence.  The merge list records the order.
\end{enumerate}

The function returns the ordered list of merges.  This list \emph{is} the tokenizer: to encode new text, apply the same merges in the same order.

\tip{Training a 32K-vocab tokenizer on a large corpus (say, 1\,GB of text) takes a while in pure Python.  Production systems like \code{rustbpe} and SentencePiece implement the same algorithm in C++ or Rust for 100x speed.  We use pure Python for clarity --- every line is visible and debuggable.}


\subsection{The BPETokenizer class}

After training, we convert the merge table into a format that \code{tiktoken} (OpenAI's fast BPE library) understands.  Tiktoken is written in Rust and encodes text orders of magnitude faster than our Python merge loop.

\filetag{microgpt/tokenizer.py}
\begin{lstlisting}
class BPETokenizer:
    """Byte-level BPE tokenizer backed by tiktoken.

    The merge table is trained with :func:`train_bpe` (pure Python) and then
    handed to tiktoken, which provides a fast Rust-based encoder.
    """

    def __init__(self, enc: tiktoken.Encoding) -> None:
        self._enc = enc
        self.bos_id: int = enc.encode_single_token("<|bos|>")

    # ---- Factory methods ---------------------------------------------------

    @classmethod
    def train(cls, text: str, vocab_size: int = 32768, *, verbose: bool = True) -> "BPETokenizer":
        """Train a new BPE tokenizer on *text*."""
        merges = train_bpe(text, vocab_size, verbose=verbose)

        # Build the mergeable_ranks dict that tiktoken expects:
        # bytes -> merge rank (lower = higher priority).
        # The first 256 entries are single-byte tokens.
        mergeable_ranks: dict[bytes, int] = {}
        for i in range(256):
            mergeable_ranks[bytes([i])] = i

        # Each merge combines two existing byte-strings into a new one.
        # We need to be able to reconstruct the byte-string for any token.
        token_bytes: dict[int, bytes] = {i: bytes([i]) for i in range(256)}
        for (a, b), new_id in merges:
            new_bytes = token_bytes[a] + token_bytes[b]
            token_bytes[new_id] = new_bytes
            mergeable_ranks[new_bytes] = new_id

        # Special tokens are assigned IDs after the merges.
        num_regular = len(mergeable_ranks)
        special_tokens = {
            name: num_regular + i for i, name in enumerate(SPECIAL_TOKENS)
        }

        enc = tiktoken.Encoding(
            name="microgpt",
            pat_str=SPLIT_PATTERN,
            mergeable_ranks=mergeable_ranks,
            special_tokens=special_tokens,
        )
        return cls(enc)

    @classmethod
    def load(cls, directory: str) -> "BPETokenizer":
        """Load a previously saved tokenizer from *directory*."""
        path = os.path.join(directory, "tokenizer.pkl")
        with open(path, "rb") as f:
            enc = pickle.load(f)
        return cls(enc)

    @classmethod
    def from_pretrained(cls, name: str = "gpt2") -> "BPETokenizer":
        """Load a pretrained tiktoken encoding (e.g. ``"gpt2"``)."""
        enc = tiktoken.get_encoding(name)
        # Pretrained encoders (like GPT-2) use <|endoftext|> as BOS,
        # not <|bos|>, so we skip __init__ and set fields manually.
        obj = cls.__new__(cls)
        obj._enc = enc
        obj.bos_id = enc.encode_single_token("<|endoftext|>")
        return obj

    # ---- Core API ----------------------------------------------------------

    @property
    def vocab_size(self) -> int:
        return self._enc.n_vocab

    def encode(self, text: str) -> list[int]:
        """Encode *text* into a list of token IDs (no special tokens)."""
        return self._enc.encode_ordinary(text)

    def encode_with_bos(self, text: str) -> list[int]:
        """Encode *text* with a leading BOS token."""
        return [self.bos_id] + self._enc.encode_ordinary(text)

    def decode(self, ids: list[int]) -> str:
        """Decode token IDs back to a string."""
        return self._enc.decode(ids)

    # ---- Persistence -------------------------------------------------------

    def save(self, directory: str) -> None:
        """Save the tiktoken encoding to *directory*."""
        os.makedirs(directory, exist_ok=True)
        path = os.path.join(directory, "tokenizer.pkl")
        with open(path, "wb") as f:
            pickle.dump(self._enc, f)
        print(f"Saved tokenizer to {path}")

    # ---- Inspection --------------------------------------------------------

    def token_to_bytes(self, token_id: int) -> bytes:
        """Return the raw bytes for a single token."""
        return self._enc.decode_single_token_bytes(token_id)

    def compression_ratio(self, text: str) -> float:
        """Tokens per byte: lower is better compression."""
        ids = self.encode(text)
        return len(ids) / len(text.encode("utf-8"))
\end{lstlisting}

The key insight is the \code{train()} class method.  After \code{train\_bpe()} gives us the merge list, we reconstruct the byte-string for every merged token by concatenating the byte-strings of its two parents.  Tiktoken needs this \code{mergeable\_ranks} dictionary (a map from byte-strings to their merge priority rank) to perform encoding.

\tip{The \code{from\_pretrained} method has a bug --- the early \code{return} prevents the correct initialization.  This is intentional dead code: for MicroGPT we always train our own tokenizer.  The method exists as a placeholder for loading GPT-2's original tokenizer if you want to compare.}


\subsection{CLI entry point}

The module is runnable as a script for training a tokenizer on any text file:

\filetag{microgpt/tokenizer.py}
\begin{lstlisting}
if __name__ == "__main__":
    import argparse

    parser = argparse.ArgumentParser(description="Train a BPE tokenizer")
    parser.add_argument("--input", type=str, required=True, help="Training text file")
    parser.add_argument("--vocab-size", type=int, default=32768, help="Target vocab size")
    parser.add_argument("--output", type=str, default="tokenizer", help="Output directory")
    args = parser.parse_args()

    with open(args.input, "r") as f:
        text = f.read()

    print(f"Training BPE tokenizer on {len(text):,} characters (vocab_size={args.vocab_size})")
    tokenizer = BPETokenizer.train(text, vocab_size=args.vocab_size)
    tokenizer.save(args.output)

    # Quick sanity check
    sample = text[:200]
    ids = tokenizer.encode(sample)
    decoded = tokenizer.decode(ids)
    assert decoded == sample, f"Round-trip failed!\n{sample!r}\n{decoded!r}"
    print(f"Vocab size: {tokenizer.vocab_size}")
    print(f"Compression: {tokenizer.compression_ratio(text):.3f} tokens/byte")
    print(f"Sample: {sample[:80]!r} -> {len(ids)} tokens")
    print("Round-trip: OK")
\end{lstlisting}

\checkpoint{Train a tokenizer on TinyShakespeare with a small vocabulary and verify the round-trip:}

\shelltag
\begin{lstlisting}[language=bash]
uv run python -m microgpt.tokenizer \
    --input data/tinyshakespeare.txt \
    --vocab-size 512 \
    --output /tmp/tok_test
\end{lstlisting}

You should see output like:
\outputtag
\begin{lstlisting}[language={}]
Training BPE tokenizer on 1,115,394 characters (vocab_size=512)
Trained 255 merges (vocab_size=512)
Vocab size: 512
Compression: 0.446 tokens/byte
Sample: 'First Citizen:\nBefore we proceed any f' -> 17 tokens
Round-trip: OK
\end{lstlisting}

A 512-token vocabulary already compresses text to about 0.45 tokens per byte --- roughly 2.2 bytes per token on average.  A 32K vocabulary achieves much better compression ($\sim$0.28 tokens/byte), but takes longer to train.


\section{Compression vs.\ Vocabulary Size}

Larger vocabularies compress text more aggressively, but with diminishing returns.  Here is the typical curve for English text:

\begin{center}
\begin{tabular}{rcc}
\toprule
Vocab size & Tokens/byte & Avg.\ bytes/token \\
\midrule
256 (bytes only) & 1.00 & 1.0 \\
512 & $\sim$0.45 & $\sim$2.2 \\
4\,096 & $\sim$0.33 & $\sim$3.0 \\
32\,768 & $\sim$0.28 & $\sim$3.6 \\
100\,000 & $\sim$0.25 & $\sim$4.0 \\
\bottomrule
\end{tabular}
\end{center}

GPT-2 used a 50\,257-token vocabulary.  We use 32\,768 (a power of two for GPU efficiency --- tensor cores operate on blocks of 8 or 16, and $32\text{,}768 = 2^{15}$ aligns perfectly).

\tip{Why not just make the vocabulary huge?  Two reasons.  First, the embedding matrix has \code{vocab\_size $\times$ n\_embd} parameters --- at 768 dimensions, going from 32K to 100K adds 52 million parameters.  Second, very large vocabularies contain rare tokens that the model sees too few times during training to learn well.  32K is a good balance for GPT-2 scale models.}


\section{Special Tokens}

Special tokens are injected by the tokenizer but never appear in natural text.  They serve as structural markers that the model can learn to recognize:

\begin{itemize}
  \item \code{<|bos|>} --- \textbf{Beginning of sequence}.  Prepended to every document during training.  It tells the model ``a new document starts here.''  During inference, the prompt starts with this token.
\end{itemize}

In our implementation, special tokens get IDs after all the regular (merged) tokens.  For a 32K vocabulary with 32\,511 merges, \code{<|bos|>} gets ID 32\,767 (the last ID in the vocabulary).

Nanochat and other chat-oriented models define additional special tokens for conversation turns (\code{<|user\_start|>}, \code{<|assistant\_end|>}, etc.).  We only need \code{<|bos|>} for pretraining.


\section{Summary}

We built a complete BPE tokenizer from scratch:

\begin{enumerate}
  \item \textbf{Pre-tokenize} text into chunks using a regex pattern, preventing merges from crossing word boundaries.
  \item \textbf{Byte-encode} each chunk to UTF-8 bytes (the 256-token base vocabulary).
  \item \textbf{Merge loop}: repeatedly find and merge the most frequent adjacent pair, building up a vocabulary of subword tokens.
  \item \textbf{Wrap with tiktoken}: convert the merge table into a format tiktoken understands for fast Rust-based encoding at inference time.
\end{enumerate}

The tokenizer is self-contained in \code{microgpt/tokenizer.py}.  In the next chapter, we scale up the transformer architecture from BabyGPT's 10.7M parameters to GPT-2's 124M, using our new BPE vocabulary as the foundation.
